% Cardinal Arithmetic
\section{Cardinal Arithmetic in the Indus System}

\subsection{Finite Cardinal Numbers}

The Indus system implements finite cardinal arithmetic with:
\begin{itemize}
    \item Domain: $\mathcal{N} = \{1, 2, 3, \ldots, n\}$ where $n \approx 12$
    \item Representation: Unary stroke notation
    \item Operations: Counting, enumeration, possibly addition
\end{itemize}

\subsection{Mathematical Structure}

\subsubsection{Monoid Structure}
The counting system forms a commutative monoid under concatenation:
\begin{align*}
    (|) \cdot (||) &= (|||) \\
    (||) \cdot (|) &= (|||)
\end{align*}

\subsubsection{Limitations}
\begin{itemize}
    \item No subtraction operator evident
    \item No multiplication operator evident
    \item No division operator evident
    \item Limited to small finite values
\end{itemize}

\subsection{Counting vs. Arithmetic}

Distinction between:
\begin{enumerate}
    \item \textbf{Enumeration}: Listing discrete quantities (present in corpus)
    \item \textbf{Calculation}: Performing arithmetic operations (not evident)
\end{enumerate}

The system supports the former but shows no clear evidence of the latter.

\subsection{Administrative Mathematics}

The mathematical capacity is sufficient for:
\begin{itemize}
    \item Inventory records (commodity counts)
    \item Transaction receipts (quantity verification)
    \item Allocation records (resource distribution)
    \item Quota specifications (required amounts)
\end{itemize}

\subsection{Comparison with Other Systems}

Similar mathematical scope:
\begin{itemize}
    \item Early Egyptian hieratic numerals (administrative contexts)
    \item Proto-cuneiform numerical tablets
    \item Chinese Shang oracle bone counts
    \item Incan quipu numerical encodings
\end{itemize}

All optimized for record-keeping rather than calculation.
